\documentclass[10pt]{scrartcl}
\usepackage[a4paper, total={5.5in, 8in}]{geometry}
\usepackage[usenames,dvipsnames,svgnames]{xcolor}
\usepackage[shortlabels]{enumitem}
\usepackage[framemethod=TikZ]{mdframed}
\usepackage{amsmath,amssymb,amsthm}
\usepackage[colorlinks]{hyperref}

\hypersetup{
	colorlinks=true,
	linkcolor=blue,
	filecolor=blue,      
	urlcolor=blue,
	pdftitle={MAT 314 Notes}
}

\usepackage{microtype}
\usepackage{mathtools}
\usepackage[headsepline]{scrlayer-scrpage}
\usepackage{thmtools}
\usepackage[textsize=small]{todonotes}

\addtolength{\textheight}{3.14cm}
\providecommand{\ol}{\overline}
\providecommand{\eps}{\varepsilon}
\providecommand{\half}{\frac{1}{2}}
\providecommand{\dang}{\measuredangle} %% Directed angle
\providecommand{\CC}{\mathbb C}
\providecommand{\FF}{\mathbb F}
\providecommand{\NN}{\mathbb N}
\providecommand{\QQ}{\mathbb Q}
\providecommand{\RR}{\mathbb R}
\providecommand{\ZZ}{\mathbb Z}
\providecommand{\dg}{^\circ}
\providecommand{\ind}[1]{\mathbf 1_{#1}}
\providecommand{\ii}{\item}
\providecommand{\alert}[1]{{\sffamily\textbf{\textcolor{magenta}{#1}}}}
\providecommand{\opname}{\operatorname}
\providecommand{\li}{\opname{li}}
\providecommand{\ls}{\opname{ls}}
\providecommand{\inv}{^{-1}}
\providecommand{\setdiff}{\smallsetminus}
\providecommand{\mb}[1]{\mathbf{#1}}
\providecommand{\mf}[1]{\mathfrak{#1}}
\providecommand{\abs}[1]{\left|{#1}\right|}
\providecommand{\dif}{\;d}

\reversemarginpar
\setlength{\marginparwidth}{3cm}

\mdfdefinestyle{mdbluebox}{roundcorner=10pt,innerbottommargin=9pt,
	linecolor=blue,backgroundcolor=TealBlue!5,}
\declaretheoremstyle[headfont=\sffamily\bfseries\color{MidnightBlue},
mdframed={style=mdbluebox},]{thmbluebox}

\mdfdefinestyle{mdredbox}{frametitlefont=\bfseries,innerbottommargin=8pt,
	nobreak=true,backgroundcolor=Salmon!5,linecolor=RawSienna,}
\declaretheoremstyle[headfont=\bfseries\color{RawSienna},
mdframed={style=mdredbox},headpunct={\\[3pt]},postheadspace=0pt,]{thmredbox}

\mdfdefinestyle{mdgreenbox}{linecolor=ForestGreen,backgroundcolor=ForestGreen!5,
	linewidth=2pt,rightline=false,leftline=true,topline=false,bottomline=false,}
\declaretheoremstyle[headfont=\bfseries\sffamily\color{ForestGreen!70!black},
mdframed={style=mdgreenbox},headpunct={ --- },]{thmgreenbox}

\mdfdefinestyle{mdorangebox}{linecolor=RedOrange,backgroundcolor=RedOrange!5,
	linewidth=2pt,rightline=false,leftline=true,topline=false,bottomline=false,}
\declaretheoremstyle[headfont=\bfseries\sffamily\color{RedOrange!70!black},
mdframed={style=mdorangebox},headpunct={ --- },]{thmorangebox}

\mdfdefinestyle{mdblackbox}{linecolor=black,backgroundcolor=RedViolet!5!gray!5,
	linewidth=3pt,rightline=false,leftline=true,topline=false,bottomline=false,}
\declaretheoremstyle[mdframed={style=mdblackbox}]{thmblackbox}


\declaretheorem[style=thmredbox,name=Problem,numberwithin=section]{problem}

\declaretheorem[style=thmredbox,name=Example,sibling=problem]{example}
\declaretheorem[style=thmredbox,name=$\bigstar$ Problem,sibling=problem]{niceprob}

\declaretheorem[style=thmbluebox,name=Theorem,sibling=problem]{theorem}
\declaretheorem[style=thmbluebox,name=Corollary,sibling=theorem]{corollary}
\declaretheorem[style=thmbluebox,name=Proposition,sibling=theorem]{prop}
\declaretheorem[style=thmbluebox,name=Lemma,sibling=theorem]{lemma}
\declaretheorem[style=thmbluebox,name=Theorem,numbered=no]{theorem*}
\declaretheorem[style=thmbluebox,name=Lemma,numbered=no]{lemma*}
\declaretheorem[style=thmgreenbox,name=Claim,numberwithin=problem]{claim}
\declaretheorem[style=thmgreenbox,name=Claim,numbered=no]{claim*}
\declaretheorem[style=thmblackbox,name=Remark,sibling=theorem]{remark}
\declaretheorem[style=thmblackbox,name=Remark,numbered=no]{remark*}
\declaretheorem[style=thmgreenbox,name=Definition,sibling=theorem]{definition}
\declaretheorem[style=thmgreenbox,name=Definition,numbered=no]{definition*}
\declaretheorem[style=thmorangebox,name=Warning,sibling=theorem]{warning}
\declaretheorem[style=thmorangebox,name=Warning,numbered=no]{warning*}
\declaretheorem[style=thmorangebox,name=Note,numbered=no]{note}
\declaretheorem[style=thmblackbox,name=Exercise,sibling=example]{exercise}
\declaretheorem[style=thmblackbox,name=Exercise,numbered=no]{exercise*}

\newenvironment{soln}{\begin{proof}[Solution]}{\end{proof}}
\newenvironment{walkthrough}{\noindent \textbf{Walkthrough.}}{\medskip}
\newlist{walk}{enumerate}{3}
\setlist[walk]{label=\bfseries (\alph*)}

\providecommand{\pow}{\operatorname{Pow}}
\providecommand{\vocab}[1]{{\textbf{\textcolor{ForestGreen}{#1}}}}
\providecommand{\lcm}{\operatorname{lcm}}
\providecommand{\ord}{\operatorname{ord}}
\providecommand{\tensor}{\otimes}

\usepackage{asymptote}
\begin{asydef}
	defaultpen(fontsize(10pt));
	unitsize(1cm);
	size(8cm); // set a reasonable default
	usepackage("amsmath");
	usepackage("amssymb");
	settings.tex="pdflatex";
	settings.outformat="pdf";
	// Replacement for olympiad+cse5 which is not standard
	import geometry;
	// recalibrate fill and filldraw for conics
	void filldraw(picture pic = currentpicture, conic g, pen fillpen=defaultpen, pen drawpen=defaultpen)
	{ filldraw(pic, (path) g, fillpen, drawpen); }
	void fill(picture pic = currentpicture, conic g, pen p=defaultpen)
	{ filldraw(pic, (path) g, p); }
	// some geometry
	pair foot(pair P, pair A, pair B) { return foot(triangle(A,B,P).VC); }
	pair orthocenter(pair A, pair B, pair C) { return orthocentercenter(A,B,C); }
	pair centroid(pair A, pair B, pair C) { return (A+B+C)/3; }
	// cse5 abbreviations
	path CP(pair P, pair A) { return circle(P, abs(A-P)); }
	path CR(pair P, real r) { return circle(P, r); }
	pair IP(path p, path q) { return intersectionpoints(p,q)[0]; }
	pair OP(path p, path q) { return intersectionpoints(p,q)[1]; }
	path Line(pair A, pair B, real a=0.6, real b=a) { return (a*(A-B)+A)--(b*(B-A)+B); }
	// cse5 more useful functions
	picture CC() {
		picture p=rotate(0)*currentpicture;
		currentpicture.erase();
		return p;
	}
	pair MP(Label s, pair A, pair B = plain.S, pen p = defaultpen) {
		Label L = s;
		L.s = "$"+s.s+"$";
		label(L, A, B, p);
		return A;
	}
	pair Drawing(Label s = "", pair A, pair B = plain.S, pen p = defaultpen) {
		dot(MP(s, A, B, p), p);
		return A;
	}
	path Drawing(path g, pen p = defaultpen, arrowbar ar = None) {
		draw(g, p, ar);
		return g;
	}
\end{asydef}

\ihead{\footnotesize\textbf{MAT 314 Spring 2025 Notes}}
\ohead{\footnotesize \today}

\title{\vspace{-2.0cm} MAT 314 Spring 2025}
\author{\large Aditya Pahuja}
\date{\large \today}
\begin{document}
\maketitle
\tableofcontents
\pagebreak

\section{Week 1: 1/28 and 1/30}

\href{https://www.math.stonybrook.edu/~rdhough/mat314_spring25/math314.html}
{Course website link}.

\subsection{The classical groups}
\begin{definition}
	Let $F$ be a ring or field.
	The \vocab{matrix group} $GL_n(F)$ is the set of
	invertible $n\times n$ matrices with entries in $F$
	under composition.
	The subgroup $SL_n(F)$ is the set of matrices with determinant $1$.
\end{definition}

\begin{prop}
	If $F$ is a ring, then $M\in GL_n(F)$ iff $\det M$ is a unit.
\end{prop}

\begin{proof}
	Recall that $\opname{adj}(M)\cdot M$ is the diagonal matrix
	whose diagonal is filled with $\det M$,
	so $\det M$ is a unit iff the diagonal matrix is invertible.
	\todo{what is the adjoint\dots?}
\end{proof}

Some groups fix a particular form:
\begin{definition}
	The \vocab{orthogonal group} $O_n(\RR)$ preserves inner product:
	\[O_n(\RR) = \{P\in GL_n(\RR) : P^tP = I_n\}.\]
\end{definition}

A variation of this is
\[O_{p,q}(\RR) = \{P\in GL_n(\RR) : P^t\}\]
\todo{text}

\begin{definition}
	The \vocab{unitary group} $U_n$ preserves a complex inner product
	$\mb x\cdot\mb y = x_1\ol{y_1} + \cdots + x_n\ol{y_n}$.
	\[U_n = \{P\in GL_n(\CC) : \}\]
\end{definition}

In particular, one can then show that the rows (and columns) of this matrix
are pairwise orthogonal. Let $P\in U_n$.
If $v$ is an eigenvector, so that $Pv = \lambda v$ for $\lambda\in\CC$, then
\[\|v\|^2 = \langle v,v\rangle = \langle Pv, Pv\rangle =
\ol v^t \ol P^t Pv = \ol v^t\ol\lambda \lambda v = |\lambda|^2\|v\|^2,\]
so $|\lambda| = 1$; therefore we can write any eigenvalue $\lambda$ as
$e^{i\theta}$.

If $v$ is an eigenvector and $w\perp v$, then
\[0 = \langle w,v\rangle = \langle Pw, Pv\rangle =
\ol\lambda\langle Pw,v\rangle,\]
so the space of vectors orthogonal to $v$ is $P$-invariant.

Note that every square matrix over $\CC$ has an eigenvalue
since $\det(P - \lambda I) = 0$ always has a solution $\lambda$
(because $\CC$ is algebraically closed).
Repeat this with orthogonal space (?) to get that any $P\in U_n$
can be diagonalized (in particular each $P\in U_n$ has an eigenbasis).

\begin{theorem}
	Every $P\in U_n$ is conjucate to a diagonal matrix of rotations;
	these are the conjugacy classes of $U_n$.
\end{theorem}

In $SU_2$, the conjugacy classes are just determined by
\[\begin{pmatrix}
	e^{i\theta} &0\\
	0&e^{-\theta}
\end{pmatrix}\]

\[SU_2 = \left\{\begin{pmatrix}
	a&b\\
	-\ol b&\ol a
\end{pmatrix} : a,b\in\CC, |a|^2 + |b|^2 = 1\right\}.\]
expanding the real and imaginary parts here gives a group law
on the $3$-sphere, analogous to the circle group for $S_1$
(an analogous result is impossible for $S_2$ because of
hairy ball theorem, or something)

\begin{definition}
	The \vocab{symplectic group} $SP_{2n}(\RR)$ preserves the
	symplectic form
	\[x_1\wedge y_1 + \cdots + x_n\wedge y_n =
	x_1\otimes y_1 - y_1\otimes x_1 + \cdots + x_n\otimes y_n - y_n\otimes x_n\]
\end{definition}

\todo{texing matrices is hard}

In general, when the name of one of these groups is prepended with an $S$,
it means we are restricting to the subgroup with det $1$
(and the adjective ``special'' is added to the front, so e.g.
$SO_n(\RR)$ is the special orthogonal group).

The matrix groups have dimensions, which are the number of
real parameters describing the group as a real manifold.
For example, $GL_n(\RR)$ has dimension $n^2$.

\begin{example}[1 dimensional group]
	The group $U_1/ SO_2(\RR)$ is the circle group,
	isomorphic to $\RR/\ZZ$, consisting of the rotation matrices
	\[\left\{\begin{pmatrix}
		\cos(2\pi\theta) & -\sin(2\pi\theta)\\
		\sin(2\pi\theta) & \cos(2\pi\theta)
	\end{pmatrix}\right\}.\]
\end{example}

\begin{example}[3 dimensional groups]
	The special linear group $SL_2(\RR)$ of $2\times 2$ matrices with det 1
	is a non-compact 3 dimensional group,
	while $SO_3(\RR)$ is a compact 3d group.
\end{example}

\begin{definition}[Quaternions]
	The group of unit quaternions is the subgroup of $SU^2$
	generated by
	\[i = \begin{pmatrix}
		i&0\\
		0&-i
	\end{pmatrix},\quad
	j = \begin{pmatrix}
		0&1\\
		-1&0
	\end{pmatrix},\quad
	k = \begin{pmatrix}
		0&i\\
		i&0
	\end{pmatrix}\]
\end{definition}

Just like $e^{i\theta} = \cos\theta + i\sin\theta$,
we have $e^{j\theta} = \cos\theta + j\sin\theta$
and $e^{k\theta} = \cos\theta + k\sin\theta$
by the power series for $e^x$.

We can compute that conjugating by $e^{i\theta}$ fixes $i$ and rotates
the $jk$-plane by $2\theta$.

unitary groups $\leftrightarrow$ unit quaternions

For a quaternion $\alpha = x_1 + x_2i + x_3j + x_4k$
The trace of a quaternion is $\opname{Tr}(\alpha) = \alpha + \ol\alpha = 2x_1$.
The norm of a quaternion is $N(\alpha) = \alpha\ol\alpha
= x_1^2 + x_2^2 + x_3^2 + x_4^2$, which is real
(and you can check this by showing that it's equal to its conjugate).

We can also show that $\ol{\alpha\beta} = \ol\beta\cdot\ol\alpha$,
so conjugation is an antihomomorphism (it flips the order of multiplication).

Recall from linear algebra that 
\[\opname{Tr}(U\inv X U) = \opname{Tr}(X),\]
i.e. similar matrices have the same trace.
This means the equator, being the set of imaginary quaternions,
is the set of matrices with trace zero.
The equator is preserved by conjugation (in the sense of groups or..?).
$SU_2$ acts on 2-sphere/equator/imaginary quaternions by conjugation.

This makes $SU_2$ a ``double cover'' of $SO_3$
in that there is a surjective homomorphism $SU_2\to SO_3$
with a kernel of size 2.

\subsection{1-parameter groups}
\begin{definition}
	The \vocab{1-parameter subgroups} of $GL_n(\RR)$ or $GL_n(\CC)$
	(which we will just denote by $GL_n$)
	are images of differentiable homomorphisms $\RR\to GL_n$.
\end{definition}

These are nice in that we can completely classify them:
\begin{theorem}
	(a) Let $A$ be an $n\times n$ real or complex matrix.
	Then, $\varphi\colon\RR\to GL_n$, given by
	$\varphi(t) = e^{tA}$, (defined in the expected way
	with the Taylor series) is a group homomorphism.
	
	(b) Conversely, if $\varphi\colon\RR\to GL_n$ is a differentiable
	group homomorphism and $\varphi'(0) = A$, then $\varphi(t) = e^{tA}$.
\end{theorem}

\begin{remark}[Convergence of $e^{tA}$]
	One can show that the series defining $e^{tA}$ converges
	under the norm $\|A\|_2 = (\sum_{i,j} A_{i,j}^2)^{1/2}$
	by proving that $\|AB\| \le \|A\| \|B\|$ and noting that
	the triangle inequality holds (because this is just the
	Euclidean norm if we treat $GL_n$ as $\RR^{n^2}$).
\end{remark}

\begin{proof}
	(a) Recall that we can prove $e^{x+y} = e^xe^y$ for $x,y\in\RR$ via
	a series argument:
	\begin{align*}
		e^{x+y} &= \sum_{n=0}^\infty \frac{(x+y)^n}{n!}\\
		&= \sum_{n=0}^\infty \frac1{n!}\sum_{k=0}^n \binom nk x^k y^{n-k}\\
		&= \sum_{n=0}^\infty \sum_{k-0}^n
		\frac{x^k}{k!}\cdot\frac{y^{n-k}}{(n-k)!}\\
		&= \sum_{j=0}^\infty \frac{x^j}{j!}\cdot
		\sum_{k=0}^\infty \frac{y^k}{k!}.
	\end{align*}
	The same argument holds with matrices because
	$rA$ and $sA$ commute for any real $r,s$. Thus
	\[\varphi(r+s) = e^{rA+sA} = e^{rA}e^{sA} = \varphi(r)\varphi(s),\]
	proving that $\varphi$ is a homomorphism,
	and differentiability is not too hard to show.
	
	(b) Conversely let $\varphi$ be a differentiable homomorphism.
	Then, $\varphi(0) = I$. Moreover, for reals $t$, $\eps$,
	\[\varphi(t + \eps) = \varphi(t)\varphi(\eps),\]
	so we can write
	\[\varphi'(t) = \lim_{\eps\to0} \frac{\varphi(t+\eps) - \varphi(t)}{\eps}
	= \lim_{\eps\to 0} \frac{\varphi(t)\varphi(\eps) - \varphi(t)}{\eps}
	= \varphi(t)\cdot\lim_{\eps\to0} \frac{\varphi(\eps) - I}{\eps}
	= \varphi(t)A.\]
	We will take for granted that the matrix exponential
	is the complete set of solutions to this differential equation.
\end{proof}

\begin{example}
	If $A = \begin{pmatrix}
		0&1\\0&0
	\end{pmatrix}$, then $A^2 = 0$, so
	\[e^{At} = I + tA = \begin{pmatrix}
		1&t\\0&1
	\end{pmatrix}.\]
\end{example}

\begin{exercise}
	What does $e^{tA}$ look like for
	\[A = \begin{pmatrix}
		0&1&0&\cdots&0\\
		0&0&1&\cdots&0\\
		\vdots&\vdots&\vdots&\ddots&\vdots\\
		0&0&0&\cdots&1\\
		0&0&0&\cdots&0
	\end{pmatrix}?\]
	(The next example uses this result.)
\end{exercise}

\begin{example}
	Recall that a Jordan block looks like a matrix $A$ from the above exercise
	plus $\lambda I$ for some real $I$. We can compute
	\[e^{\lambda I + A} = e^{\lambda I} e^A = e^\lambda I\begin{pmatrix}
		1&t&\cdots&\frac{t^n}{n!}\\
		\vdots&\vdots&\ddots&\vdots\\
		0&0&\cdots&t\\
		0&0&\cdots&1
	\end{pmatrix}\]
	using the result of the exercise.
	We can also show that $e^{B\inv AB} = B\inv e^AB$,
	so this should let us compute the exponential of an $A\in GL_n(\CC)$
	pretty easily maybe?
\end{example}

\begin{example}
	If $A$ is the $\pi/2$ rotation matrix ($j$ from the quaternions)
	then $e^{tA}$ is the $t$-radian rotation matrix.
\end{example}

\begin{prop}
	(a) If $A^t = -A$ for a real matrix $A$ then $e^A$ is orthogonal,
	and if $\ol{A^t} = -A$ for a complex matrix $A$ then $e^A$ is unitary.
	
	The other two parts here say that the converses hold:
	
	(b) The 1-parameter subgroups of $O_n$ are $t\mapsto e^{tA}$
	for a skew-symmetric $A$.
	
	(c) The 1-parameter subgroups of $U_n$ are $t\mapsto e^{tA}$
	for a skew-Hermitian $A$.
\end{prop}

\begin{proof}
	We just show the parts about complex matrices; the results for real matrices
	follow because Hermitian and symmetric mean the same thing for reals
	and unitary and orthogonal mean the same thing for reals.
	
	If $\ol{A^t} = -A$, then we can show by expanding the series that
	\[\ol{e^{A^t}} = e^{\ol{A^t}} = e^{-A}\]
	so that
	\[(\ol{e^{A}})^t e^A = I\]
	which implies $e^A$ s unitary. For the other direction,
	if $\varphi\colon\RR\to GL_n(\CC)$ is a differentiable homomorphism,
	\[\frac{d}{dt}(\ol{\varphi(t)})\varphi(t) = \ol{\varphi'(t)^t}\varphi(t)
	+ \ol{\varphi(t)^t\varphi'(t)}\]
	(after checking that the product rule holds, which we can do by
	just expanding the matrix product).
	This tells us that
	\[\ol{\varphi'(0)}^t + \varphi'(0) = 0\]
	so that $A := \varphi'(0)$ is skew-Hermitian, and we already know that
	$\varphi$ has to look like $e^{t\varphi'(0)}$ from
	the structure theorem we proved about one-parameter groups.
\end{proof}

\begin{lemma}
	$e^{\opname{tr}A} = \det e^A$.
\end{lemma}

\begin{proof}
	We know that $A$ is similar to $C = B\inv AB$,
	a direct sum of Jordan blocks, and the determinant
	and trace are preserved under similarity,
	so it suffices to check this for a Jordan block;
	but this is just a simple computation,
	looking at the example from earlier.
\end{proof}

\begin{prop}
	The one-parameter subgroups of $SL_n(\RR)$ are $e^{tA}$
	where $A$ is a real matrix with a trace of zero.
\end{prop}

\begin{proof}
	$e^{tA}$ has determinant $1$ if and only if the trace of $tA$ is zero.
\end{proof}

\subsection{Lie algebras}

\begin{definition}
	The \vocab{Lie algebra} of a matrix group is the set of
	tangent vectors to the group at the identity.
\end{definition}

Concretely, these tangent vectors manifest as derivatives at zero
of differentiable functions (not necessarily homomorphisms)
$\varphi\colon\RR\to G$ such that $\varphi(0) = I$.

\begin{prop}
	The Lie algebra of $O_n$ is skew-symmetric real matrices.
\end{prop}

\begin{proof}
	The one-parameter subgroup characterization above shows that
	skew-symmetric matrices are all tangent vectors.
	In the other direction, we have that if $\varphi\colon\RR\to O_n$
	is differentiable then
	\[\varphi(t)^t\varphi(t) = I
	\implies \frac{d}{dt}[\varphi(t)^t \varphi(t)]_{t=0} = 0,\]
	and expanding with the product rule and using $\varphi(0) = I$
	shows that $\varphi'(0)$ is skew-symmetric.
\end{proof}

\begin{lemma}
	If $\varphi\colon\RR\to GL_n$ is a differentiable function,
	$\varphi(0) = I$, and $\varphi'(0) = A$,
	then $\frac{d}{dt} [\det\varphi]_{t=0} = \opname{tr}A$.
\end{lemma}

\begin{proof}[Proof Sketch]
	Write $\varphi(t) = \left(\varphi_{ij}(t)\right)_{i,j}$
	and then expand $\det\varphi(t)$ in terms of these coordinate functions;
	$\varphi_{i,\sigma(i)}(0)$ is equal to $1$ if $i = \sigma(i)$
	and zero otherwise.
\end{proof}

\begin{prop}
	The Lie algebra of $SL_n(\RR)$ is real matrices with trace zero.
\end{prop}

\begin{proof}
	Use the one-parameter group characterization and the preceding lemma.
\end{proof}

\begin{definition}
	The \vocab{bracket} of a Lie algebra is $[A,B] = AB - BA$.
\end{definition}

The bracket maps pairs of elements of the lie algebra
to an element of the Lie algebra and satisfies the \vocab{Jacobi identity},
\[[A,[B,C]] + [B,[C,A]] + [C,[A,B]] = 0\]
(which is sort of a proxy for associativity); this is easy to show
by just expanding (all orderings of $A$, $B$, $C$ appear twice
with different signs).

\begin{definition}
	More abstractly, a Lie algebra in its full generality is
	a real vector space $V$ with a bracket
	$[\bullet,\bullet]\colon V\times V\to V$.
	The bracket satisfies the following properties:
	\begin{enumerate}[(1)]
		\ii Bilinearity,
		\[[cv_1 + v_2, w] = c[v_1, w] + [v_2, w],\qquad
		[v, cw_1 + w_2] = c[v, w_1] + [v, w_2].\]
		\ii Skew symmetry,
		\[[v,w] = -[w,v]\]
		(so in particular $[v,w] = 0$).
		\ii Jacobi identity,
		\[[u,[v,w]] + [v,[w,u]] + [w,[u,v]] = 0.\]
	\end{enumerate}
\end{definition}

\begin{remark}
	Baker-Campbell-Hausdorff identity --- non-convergent infinite expansion
	$e^Ae^B = e^C$ for $C$ in terms of nested brackets
\end{remark}

\begin{definition}
	Define \vocab{left translation} by a matirx $P$
	by $m_P\colon G\to G$, $x\mapsto Px$.
	(This is just the standard left multiplication group action
	of $GL_n$ on $G$).
\end{definition}

\begin{definition}
	A \vocab{manifold of dimension $d$} is a topological space
	in which every point has a neighborhood homeomorphic to
	(an open ball in) $\RR^d$.
\end{definition}

\begin{lemma}
	The map $A\mapsto e^A$ is a homeomorphism from
	a neighborhood of $0\in\RR^{n\times n}$ to a neighborhood of
	$I\in GL_n(\RR)$.
\end{lemma}

\begin{proof}
	It's easy to show that the function is continuous.
	If $\|B\| < 1$ (where the norm is the one defined earlier,
	essentially just the Euclidean norm), then the series
	\[\log(I + B) = B - \frac{B^2}{2} + \frac{B^3}{3} - \frac{B^4}{4} +\cdots\]
	is convergent. Then,
	\[\|e^A - I\| \le \sum_{k=1}^\infty \frac{\|A^k\|}{k!}
	\le \sum_{k=1}^\infty \frac{\|A\|^k}{k!} = e^{\|A\|}-1,\]
	so the image of a neighborhood is a neighborhood,
	and if we just choose $\|A\|$ small enough that $\log(e^A)$ converges
	then we get a continuous inverse on the neighborhood.
\end{proof}

\begin{prop}
	$O_n$ is a manifold of dimension $1 + 2 + \cdots + (n-1)$.
\end{prop}

\pagebreak
\section{Week 3: 2/4 and 2/6}
\subsection{Normal subgroups of matrix groups}
Recall the following definition:
\begin{definition}
	A space $X$ with a topology $\mathcal T$ is \vocab{path-connected}
	if, given $x,y\in X$, there is a continuous function $\varphi\colon[0,1]\to X$
	such that $\varphi(0) = x$ and $\varphi(1) = y$.
\end{definition}

This also implies connectedness in the standard topological sense:
\begin{prop}
	A path-connected space $X$ is also connected,
	meaning it is not the disjoint union of two non-empty open sets.
\end{prop}
\begin{proof}
	If $X = U\sqcup V$ for non-empty open $U,V\subseteq X$,
	take $x\in U$ and $y\in V$, and then take a path $\varphi$
	from $x$ to $y$. Then the image of $\varphi$ is connected
	because $[0,1]$ is connected, but this clashes with the fact that
	$\varphi([0,1])\cap U$ and $\varphi([0,1])\cap V$ are disjoint
	non-empty open in the subspace $\varphi([0,1])$.
\end{proof}

\begin{prop}
	Let $G$ be a path-connected matrix group.
	Let $H < G$ contain a non-empty open $U\subseteq G$.
	Then, $H = G$.
\end{prop}

\begin{proof}
	Recall that if $U$ is open then translation in the group is continuous.
	Moreover,
	\[G = \bigcup_{g\in G} gU.\]
	If $C$ is a coset of $H$, then
	\[C = \bigcup_{g\in gU\subseteq C} gU.\]
	If there were more than one coset, the cosets are pairwise disjoint
	and, by the expression above, open, contradicting the fact that $G$
	is path-connected.
\end{proof}

\begin{theorem}
	The only proper normal subgroup of $SU_2$ is $\{\pm I\}$,
	and $SO(3)$ is simple.
\end{theorem}

\begin{proof}
	The second part follows from the first because of the double cover thing.
	
	$SU_2$ is path-connected because it's topologically the same as $S^3$.
	\todo{review}
	Suppose $N\trianglelefteq SU_2$ and let $P\in N$ not be $\pm I$.
	Then $N$ contains the conjugacy class of $P$.
	We showed a while ago that each conjugacy class in $SU_2$
	is represented by exactly one matrix of the form
	\[\begin{pmatrix}
		e^{i\theta} & 0\\
		0 & e^{-i\theta}
	\end{pmatrix},\]
	so $C$, the conjugacy class of $P$, contains all matrices of a given trace
	$t\in[-2,2]$.
	Choose $P_1\ne P_2\in C$. Then $N$ contains $P_1\inv P_1 = I$
	and $P_1\inv P_2 = Q$ ($\ne I$?).
	The conjugacy class contains a path connecting $P_1$ to $P_2$,
	so $N$ contains a path connecting $I$ to $Q$.
	Now the trace of $Q$ is less thna $2$ and the trace of $I$ is $2$,
	so $N$ contains all matrices of trace between $\opname{Tr}(Q)$ and $2$,
	which implies $N$ contains an open set about the origin,
	implying $N = SU_2$.\todo{what?}
\end{proof}

\begin{theorem}
	Let $F$ be a field with finitely many elements such that $|F| > 5$.
	\begin{enumerate}[(a)]
		\ii $SL_2(F)$ has $\{\pm I\}$ as its only normal subgroup.
		\ii $PSL_2(F) = SL_2(F)/\{\pm I\}$ is a simple group.
	\end{enumerate}
\end{theorem}

\begin{proof}
	Obviously the two parts of the theorem are equivalent.
	
	\begin{claim}
		If $F$ is a finite field with $|F| > 5$, then
		there is $r\in F$ such that $r^2 = s\notin\{0, \pm 1\}$.
	\end{claim}
	\begin{proof}
		The map $x\mapsto x^2$ maps $0$ to $0$ and collapses pairs
		of elements of $F$ onto one element; if $|F| > 5$,
		then the image is therefore just too big.
		\todo{this feels off for parity reasons?}
	\end{proof}

	Let $V = F^2$ be a $2$-D vector space over $F$ and let
	$N\trianglelefteq SL_2(F)$. Moreover, let $A\ne I$ be in $N$.
	Also, let $r^2 = s\notin\{0,\pm1\}$. We will construct a matrix
	\[B = A(PA\inv P\inv)\in N\]
	with eigenvalues $s$ and $\frac1s$.
	
	Let $v_1\in V$ be not an eigenvector (possible because $A$ is not
	a multiple of $I$?). Then if $v_2 = Av_1$, $v_1,v_2$ form a basis of $V$.
	Now define
	\[P = \begin{pmatrix}
		v_1&v_2
	\end{pmatrix}\begin{pmatrix}
		r&0\\
		0&1/r
	\end{pmatrix}\begin{pmatrix}
		v_1&v_2
	\end{pmatrix}\inv,\]
	then we can compute that the eigenvalues are indeed $s$ and $\frac1s$.
	
	\begin{claim}
		The matrices in $SL_2(F)$ with $s$, $\frac1s$ as eigenvalues
		form a conjugacy class.
	\end{claim}
	\begin{proof}
		If $M$ has eigenvalues $s$, $\frac1s$, then in $GL_2(F)$
		it is diagonalizable via a matrix $L$ (i.e. $L\inv ML$
		is diagonal with eigenvalues along diagonal);
		we can multiply $L$ by $\begin{pmatrix}
			x&0\\0&1
		\end{pmatrix}$ to force determinant to be $1$,
		meaning the diagonalized matrix is in $N$.
	\end{proof}
	
	$N$ thus contains
	$\begin{pmatrix}
		1&x\\0&1
	\end{pmatrix}$ and its transpose for arbitrary $x$, since
	\[\begin{pmatrix}
		\frac1s & 0\\0 & s
	\end{pmatrix}
	\begin{pmatrix}
		s & sx\\0 & \frac1s
	\end{pmatrix} =\begin{pmatrix}
		1&x\\0&1
	\end{pmatrix}\]
	and similar computation for the transpose.
	
	Now we show that
	\[\left\{\begin{pmatrix}
		1&x\\0&1
	\end{pmatrix},
	\begin{pmatrix}
		1&x\\0&1
	\end{pmatrix} : x,y\in F\right\}\]
	generates $SL_2(F)$. The point is that these two matrices
	let you do the ``add multiple of one row to another''
	row operation, and that's enough to get all of $SL_2$.
	Yay!\dots?
\end{proof}

\subsection{Representations of finite groups}
\begin{definition}
	Let $GL_n = GL_n(\CC)$. A \vocab{matrix representation}
	of a finite group $G$ is a group homomorphism $R\colon G\to GL_n$
	such that $R_g\cdot R_h = R_{gh}$.
	The \vocab{dimension} of the representation is $n$.
\end{definition}

\begin{example}
	The diherdal group $D_3$ is generated by $\frac{2\pi}3$ rotation, $a$,
	and by reflection in the $x$-axis, $b$, so
	\[D_3 = \{I, a, a^2, b, ab, a^2b\},\]
	giving an obvious $2$-dimensional representation by setting
	$R_a$ and $R_b$ to be the corresponding $2\times 2$ matrices.
\end{example}

\begin{example}
	The symmetric group $S_n$ has an obvious $n$-dimensional representation
	via permutation matrices, called the permutation representation.
	It also has a $1$-dimensional representation via
	$\sigma\mapsto\opname{sgn}(\sigma)$;
	this can be computed as the determinant of the corresponding
	permutation matrix (because switching rows flips sign of determinant).
\end{example}

\begin{definition}
	Given a matrix representation $R\colon G\to GL_n$,
	the \vocab{character} $\chi$ is defined as $\chi(g) = \opname{Tr} R_g$.
	This character is constant on conjugacy classes.
\end{definition}

Slightly generalizang the previous definition:
\begin{definition}
	A \vocab{representation} of a group $G$ in a complex vector space
	$V$ is a group homomorphism $\rho\colon G\to GL(V)$
	satisfying $\rho_g\cdot \rho_h = \rho_{gh}$.
\end{definition}
Of course, anchoring yourself to a basis gives a matrix representation.
This should be thought of as the natural action of $G$ on $V$.

\begin{definition}
	Given representations $\rho_1\colon G\to GL(V_1)$ and
	$\rho_2\colon G\to GL(V_2)$ and an invertible linear map
	$T\colon V_1\to V_2$, we say that $\rho_1$ and $\rho_2$
	are \vocab{equivalent} representations if
	\[T(gv) = gT(v),\]
	abbreviating $\rho(g)\cdot v$ as simply $gv$
	to align with the intuition of group actions.
\end{definition}

\begin{definition}
	A vector $v\in V$ is $G$-invariant if $gv = v$ for all $g\in G$,
	and a subspace $W\subseteq V$ is $G$-invariant if
	$gW\subseteq W$ for all $g\in G$.
\end{definition}

\begin{example}
	It is apparently an important fact that the vector
	\[\tilde{v} = \frac{1}{|G|}\sum_{g\in G} gv\]
	is always $G$-invariant.
\end{example}

The idea of $G$-invariant subspaces is important because it means that,
given a representation $\rho\colon G\to GL(V)$,
$\rho$ restricts to a representation $W$.

\begin{lemma}
	Let $W$ be $G$-invariant. Then $gW = W$ for all $g$.
\end{lemma}

\begin{proof}
	Elements of $GL(V)$ are invertible by definition,
	so they preserve dimension, hence $gW$ is a subspace of $W$
	whose dimension matches that of $W$.
\end{proof}

Recall that we can combine two vector spaces by taking their direct sum:
specifically $V = W_1\oplus W_2$ if $W_1\cap W_2 = \{0\}$ and
$W_1 + W_2 = V$; in this case each $v\in V$ can uniquely be written as
$w_1 + w_2$ for some $w_1\in W_1$, $w_2\in W_2$.
This extends to vector spaces:

\begin{definition}
	Let $\rho\colon G\to GL(V)$ be a group representation
	and suppose that $V = W_1\oplus W_2$ where $W_1$, $W_2$ are $G$-invariant.
	Then $\alpha = \rho|_{W_1}$ and $\beta = \rho|_{W_2}$ are representations
	and we write $\rho = \alpha\oplus\beta$.
\end{definition}

If you choose a basis of $V$ which is made of bases of $W_1$ and $W_2$
then the matrix of $\rho$ has form
\[R_g = \begin{pmatrix}
	A_g & 0\\
	0 & B_g
\end{pmatrix},\]
and in this case we say $\rho$ is reducible.
\todo{this is same as indecomposable here ?}

\begin{definition}
	A representation $\rho\colon G\to GL(V)$ is \vocab{irreducible}
	if the only $G$-invariant subspaces of $V$ are $0$ and $V$.
\end{definition}

\subsection{Unitary representations}

\begin{definition}
	$V$ is a \vocab{Hermitian} vector space if it is a vector space over $\CC$
	carrying a \vocab{sesquilinear form}, meaning a function
	$\langle\bullet,\bullet\rangle\colon V\times V\to\CC$ such that
	\begin{itemize}
		\ii $\langle av_1 + v_2,w\rangle = a\langle v_1,w\rangle +
		\langle v_r,w\rangle$
		\ii $\langle v,w\rangle = \ol{\langle w,v\rangle}$
		\ii $\langle v,v\rangle\ge 0$ with equality iff $v=0$.
	\end{itemize}
	We write $\|v\|^2 = \langle v,v\rangle$.
	A linear operator on $V$ is \vocab{unitary} if
	$\langle Tv,TW\rangle = \langle v,w\rangle$.
\end{definition}

Given an orthogonal basis of a Hermitian space, the matrix of
a unitary operator satisfies $A^* = A\inv$ where
$A^*$ is the transpose of the conjugate.

\begin{definition}
	If $V$ is Hermitian and $\rho\colon G\to GL(V)$ a representation of $G$,
	then we say that $\rho$ is \vocab{unitary} if $\rho(g)$
	is a unitary operator on $V$ for all $g\in G$, meaning
	$\langle gv,gw\rangle = \langle v,w\rangle$ for all $g\in G$, $v,w\in V$.
	A matrix representation $R$ is unitary if $R_g\in U_n$
	is a unitary matrix for all $g\in G$.
\end{definition}

This following lemma is important in addressing a subtlety of
how we defined reducibility/irreducibility:
we don't know as of yet that having a nontrivial $G$-invariant subspace
is the same thing as having a decomposition of $V$ as a direct sum
of $G$-invariant subspaces.
\begin{lemma}
	If $V$ is a Hermitian space, $\rho\colon G\to GL(V)$ is a
	unitary representation, and $W<V$ is $G$-invariant,
	then $W^\perp$, the set of vectors in $V$ orthogonal to $W$,
	is also $G$-invariant.
\end{lemma}

\begin{proof}
	If $w\in W$ and $u\in W^\perp$,
	\[\langle u,w\rangle = 0 = \langle gu, gw\rangle.\]
	Since this holds for all $w\in W$ because $\rho(g)\colon W\to W$ is onto,
	$gu\in W^\perp$ as well.
\end{proof}

\begin{corollary}
	A unitary representation $\rho$ on a Hermitian space $V$
	is the direct sum of irreducible representations.
\end{corollary}

\begin{proof}
	If $\rho$ is irreducible then we're done.
	
	Otherwise, let $W$ be an irreducible subrepresentation of $\rho$.
	We then get $V = W\oplus W^\perp$ by the lemma,
	and we can apply induction on $W^\perp$ since its dimension
	is strictly smaller than that of $V$.
\end{proof}

The big thing, though, is that ``unitary representation on a Hermitian space''
isn't actually that big of a request to make!
\begin{theorem}
	Let $\rho\colon G\to GL(V)$ be a representation if a complex
	vector space $V$. Then there is an inner product on $V$
	(making $V$ Hermitian) such that $\rho$ is unitary
	with respect to this inner product.
\end{theorem}

\begin{proof}
	First we need an inner product to work with.
	Fix $v_1$, \dots, $v_n$ to be a basis. Then we can just sort of steal
	the $\CC^n$ inner product:
	\[\langle x_1v_1+\cdots+x_nv_n,y_1v_1+\cdots+y_nv_n\rangle
	= x_1\ol{y_1} + \cdots + x_n\ol{y_n}.\]
	To make $\rho$ unitary, we do an averaging trick.
	Define a new inner product by
	\[\langle v,w\rangle_1 = \frac{1}{|G|}\sum_{g\in G}\langle gv,gw\rangle.\]
	(The sum behaves well with sesquilinearity so this is indeed
	an inner product).
	Now $\rho$ is unitary with respect to this inner product
	because multiplying by $G$ will just permute the things
	inside the summation!
\end{proof}

Collating everything:
\begin{corollary}[Maschke's theorem]
	Let $G$ be a finite group and $\rho\colon G\to GL(V)$ a representation
	over a complex vector space $V$.
	\begin{enumerate}[(a)]
		\item $\rho$ can be expressed as a direct sum
		of irreducible subrepresentations.
		\ii There is a basis of $V$ that makes the matrix representation
		unitary.
		\ii If $R\colon G\to GL_n$ is a matrix representation,
		there is an invertible $P\in GL_n$ such that $P\inv R_g P$
		is unitary for all $g\in G$.
		\ii Every finite subgroup of $GL_n$ is conjugate to
		a subgroup of $U_n$.
	\end{enumerate}
\end{corollary}

\begin{corollary}
	Let $A\in GL_n$ and suppose that $A^k = I$ for some $k$.
	Then, $A$ is diagonalizable.
\end{corollary}

\begin{proof}
	The group generated by $A$ is a finite (cyclic) group,
	so is conjugate to a subgroup of $U_n$, and unitary matrices
	are all diagonalizable.
\end{proof}

\subsection{Characters}
\begin{definition}
	Let $\rho\colon G\to GL(V)$ be a representation.
	The \vocab{character} of $\rho$ is $\chi\colon G\to\CC$, defined by
	$\chi(g) = \opname{Tr}\rho(g)$; the trace can be computed
	in any basis. We say $\chi$ is irreducible iff $\rho$ is.
\end{definition}

\begin{prop}
	Let $\chi$ be the character of any representation $\rho$.
	\begin{enumerate}[(a)]
		\ii $\chi(1) = \dim\rho$.
		\ii $\chi(h\inv gh) = \chi(g)$.
		\ii If $\ord(g) = k$, then $\chi$ is the sum of $\dim\rho$ powers
		of $\zeta_k = e^{2\pi i/k}$.
		\ii $\chi(g\inv) = \ol{\chi(g)}$.
		\ii $\rho\oplus\rho'$ has $\chi + \chi'$ as its character.
		\ii If $\rho\cong\rho'$, then $\chi = \chi'$.
	\end{enumerate}
\end{prop}

\begin{proof}
	(a) The matrix of $R_1$ is just $I_d$ in any basis
	since $\rho(1)$ is the identity.
	
	(b) $\opname{Tr}(h\inv gh) = \opname{Tr}(ghh\inv) = \opname{Tr}(g)$.
	
	(c) dont understand this one
	
	(d) Make the matrix $R_g$ unitary and then use the fact that
	its inverse is the conjugate transpose (and the transpose
	doesn't touch the trace).
	
	(e) Choose a basis so that the matrix looks like
	two blocks; then the two blocks each contribute
	their own trace, meaning the character of the direct sum
	is just the sum of the traces.
	
	(f) 
\end{proof}

\begin{definition}
	There is a Hermitian inner product on $\mathcal H$, the space of
	\vocab{class functions} $G\to\CC$ that are invariant under conjugation.
	Specifically
	\[\langle f_1,f_2\rangle = \frac{1}{|G|}\sum_{g\in G}\ol{f_1(g)}f_2(g)\]
	(note that the \emph{first} term in the product is conjugated,
	not the second).
\end{definition}

If $C_1$, \dots, $C_r$ are the conjugacy classes and $c_i\in C_i$ are chosen,
\[\langle f_1, f_2 = \frac{1}{|G|}\sum_{i=1}^r |C_i|\ol{f_1(c_i)}f_2(c_i)\]
because the $f_i$ are constant on conjugacy classes.
Thus we can identify $\mathcal H$ with $\CC^r$
with the inner product scaled by $|C_i|/|G|$ (so just a weighted sum).

The big theorem is the following:
\begin{theorem}
	Let $G$ be a finite group.
	\begin{enumerate}[(a)]
		\ii The irreducible characters are an
		orthonormal basis for $\mathcal H$.
		\ii In particular, the number of irreducible representations of $G$
		is equal to the number of conjugacy classes (up to equivalence).
		\ii If $\rho_1$, \dots, $\rho_r$ are the irreducible representations
		with $d_i = \dim\rho_i$, then $d_i\mid |G|$ and
		$d_1^2 + \cdots + d_r^2 = |G|$.
	\end{enumerate}
\end{theorem}

\begin{corollary}
	Two representations are isomorphic if and only if
	their characters are equal. (WTF)
\end{corollary}

\begin{proof}
	Write $\rho = \bigoplus_{i=1}^r n_i\rho_i$ and
	$\rho'$ similarly. Then being isomorphic follows from only having
	$r$ characters (one per conjugacy class).
\end{proof}

\begin{corollary}
	If $\chi$, $\chi'$ are character,s their inner product is
	a positive integer.
\end{corollary}

\begin{proof}
	It's just gonna be $n_1n'_1 + \cdots + n_rn'_r$.
\end{proof}

If $r$ is the number of conjugacy classes, a \vocab{character table}
of $G$ is the $r\times r$ matrix listing the character value
of each irreducible representation at each conjugacy class.
Scaling properly, this is an orthogonal (orthonormal?) matrix.

\begin{definition}
	A 1-dimensional character of a group $G$ is a
	homomorphism of $G$ in a 1-dimensional vector space,
	i.e. a homomorphism $G\to C^\times$.
\end{definition}

This means $\chi(g) = \rho_g = R_g$ and $\chi(g\cdot h) = \chi(g)\chi(h)$.
The image of $\chi$ is in $\CC^\times$, which is abelian, so
$\chi$ maps commutators to $1$. Thus $\chi$ descends to a homomorphism
from the abelianization $G^{\mathrm{ab}} = G/[G,G]\to \CC^\times$.

An important special case: if $G$ is a finite abelian group,
then all irreducible representations are $1$-dimensional
(the underlying algebra is commutative?)
and the number of such representations is the size of $|G|$
(because that's the number of conjugacy classes).
These are all $1$-dimensional characters.

\section{Week 3: 2/11 and 2/13}
We can make a representation for any arbitrary finite group
in analogy with the permutation representation for $S_n$.
If $G$ acts on a set $S$, then $s\mapsto g\cdot s$ is bijective\todo{wait why},
so we can associate a permutation $\sigma_g$ of $S$ to each $g\in G$.
Then we can make a $|S|$-dimensional vector space with standard basis
$\{e_s\}_{s\in S}$ such that $R_ge_s = e_{g\cdot s}$,
which gives us a representation.

\begin{lemma}
	The character $\chi(g)$ of this permutation representation is
	\[\chi(g) = \#\{s\in S< g\cdot s = s\},\]
	the number of fixed points of the action of $g$ on $S$.
\end{lemma}

This implies the following:
\begin{lemma}
	Let $R$ be  a permutation representation of $G$ acting on a finite set $S$.
	Also, let $\chi_1$ be the character of the trivial representation.
	Then, $\langle x, x_1\rangle > 0$ always, so $R$ contains
	the trivial representation as a subrepresentation.
\end{lemma}
\begin{proof}
	A direct computation:
	\begin{align*}
		\langle \chi,\chi_1\rangle = \frac{1}{|G|} = \sum_{g\in G}\ol{\chi(g)}\chi_1(g) \ge \frac{|S|}{|G|} > 0
	\end{align*}
	since $\chi(e)$ contributes $|S|$ to the sum and $\chi_1$ is always $1$.
\end{proof}
\todo{estoy confused}

\begin{definition}
	The regular representation of a group $G$ is the permutation representation
	when $G$ acts on itself by left multiplication.
\end{definition}
The character of the regular representation is
\[\chi_{reg}(g) = \begin{cases}
	|G| & g=e\\
	0 & g\ne e
\end{cases}\]
noting that $g\cdot h\ne h$ if $g\ne e$ (so no fixed points).
The dimension of $\rho_{reg}$ is $|G|$.

For any character $\chi$,
\[\langle\chi_{reg}, \chi\rangle = \frac{1}{|G|}\sum_{g\in G}\ol{\chi_{reg}(g)}\chi(g) = \chi(e) = \dim\rho\]

\begin{corollary}
	If $\chi_1$, \dots, $\chi_r$ are the characters of the irreducible
	representations of $G$ then $\chi_{reg} = d_1\chi_1 + \cdots + d_r\chi_r$
	and the regular representation is the corresponding
	direct sum of the irreducible subrepresentations.
	This exactly gives us $|G| = d_1^2 + \cdots + d_r^2$.
\end{corollary}

\begin{definition}
	Let $\rho\colon G\to GL(V)$ and $\rho'\colon G\to GL(V)$ be representations.
	A linear transformation $T\colon V\to V'$ is $G$-invariant if
	\[Tgv = gTv\]
	for all $v\in V$, $g\in G$.
	Equivalently
	\[\rho'(g)\inv T\rho(g) = T\]
	or that [diagram] commutes.
\end{definition}

In this language a bijective $G$-invariant transformation is
an isomorphism of representations.

In matrix form, after choosing bases $B$, $B'$ for $V$, $V'$,
if $M$ is the matrix for $T$, then
\[MR_g = R'_gM\]
is the same thing as $G$-invariance. In particular $R_g$ and $R'_g$
are similar matrices.

\begin{lemma}
	The kernel and image of a $G$-invariant transformation are
	$G$-invariant subspaces.
\end{lemma}

\begin{proof}
	Suppose $T\colon V\to V'$ is $G$-invariant.
	If $Tv = 0$ then $Tgv = gTv = 0$, so $gv\in\ker T$ implies $G$-invariance.
	If $Tv = w$, $Tgv = g\cdot w$, which implies that
	if $w$ is in the image of $T$ then $gw$ is also in the image of $T$,
	so the image is also $G$-invariant.
\end{proof}

\begin{lemma}[Schur]
	(a) Let $\rho$, $\rho'$ be irreducible representations of $G$
	over spaces $V$, $V'$ respectively. If $T\colon V\to V'$
	is $G$-invariant then $T$ is either an isomorphism or the zero map.s
	
	(b) If $\rho$ is an irreducible representation of $G$ over $V$,
	then $T\colon V\to V$ being $G$-invariant implies that
	$T = c\cdot I$ is a map multiplying every vector by a scalar.
\end{lemma}

\begin{proof}
	If $T$ is $G$-invariant then its kernel is a subrepresentation of $\rho$,
	so the kernel is either zero or $V$. In the first case,
	we know that the image is a nonzero subrepresentation of $\rho'$
	(because $T$ isn't the zero map) so the image is $V'$,
	implying bijection hence isomorphism.
	The second case trivially implies $T$ is the zero map.
	
	(b) Let $\lambda\in\CC$ be an eigenvalue of $T$.
	Then $T-\lambda I$ is $G$-invariant and has nontrivial kernel,
	which implies that $T-\lambda I$ is the zero map and so $T = \lambda I$.
\end{proof}

Given any linear map $T\colon V\to V'$, form the $G$-average
\[\tilde T = \frac1{|G|}\sum_{g\in G}\rho'(g)\inv T\rho(g)\]
We can compute directly that this map is $G$-invariant. \todo{write out}
We define $\tilde M$ analogously for a matrix $M$
and matrix representations $R$, $R'$.

\begin{prop}
	Let $\rho$ be an irreducible representation of $G$ in $V$
	and $T\colon V\to V$ a linear map. Then $T$ and $\tilde T$
	have the same trace. If the trace of $T$ is nonzero then $\tilde T$
	is not the zero map.
\end{prop}

\begin{proof}
	The addends in the expression for $\tilde T$ are conjugates of $T$,
	so they all have the same trace as $T$.
\end{proof}

\subsection{Orthogonality}
Let $\mathcal M = \CC^{m\times n}$ be the $m\times n$ matrices
with entries in $\CC$. Consider the map $F\colon \mathcal M\to\mathcal M$
given by
\[F(M) = AMB\]
for $A\in\CC^{m\times m}$, $B\in\CC^{n\times n}$.
Clearly $F$ is a linear map.

\begin{prop}
	The trace of $F$ is $(\opname{Tr}A)\cdot (\opname{Tr}B)$.
\end{prop}

\begin{proof}
	The standard basis vectors in $\CC^{m\times n}$ are the matrices
	whose entries are all zero except for a single $1$, say
	row $i$, column $j$; we will call the basis element $e_{i,j}$. Then
	\[\opname{Tr}F = \sum_{i,j}\langle F(e_{i,j}) e_{i,j}\rangle.\]
	Each addend is the $i,j$ entry in $Ae_{i,j}B$
\end{proof}

\todo{what the fuck happened for the second half of Tuesday's lecture}

\subsection{Tensor interlude: tensor products of vec spaces and reps}
\begin{definition}[Tensor product of vector spaces]
	The \vocab{tensor product} of two vector spaces $V_1$, $V_2$,
	denoted $V_1\otimes V_2$, is the span of
	\[\{v_1\otimes v_2 : v_1\in V_1,v_2\in V_2\}\]
	where the product is multi-linear; namely,
	\begin{align*}
		(cv_1 + v_1')\otimes v_2 &= c(v_1\otimes v_2) + v_1'\otimes v_2\\
		v_1\otimes (cv_2 + v_2') &= c(v_1\otimes v_2) + v_1\otimes v_2'.
	\end{align*}
	In particular, if $e_1$, \dots, $e_m$ is a basis for $V_1$
	and $f_1$, \dots, $f_n$ is a basis for $V_2$, then the set of
	$e_i\oplus f_j$ is a basis for $V_2\otimes V_2$.
\end{definition}

We can then tensor together representations:
\begin{definition}[Tensor product of representations]
	Given representations $\rho_1\colon G\to GL(V_1)$
	and $\rho_2\colon G\to GL(V_2)$ of $G$, the tensor product of
	the representations is $\rho_1\otimes \rho_2\colon G\to GL(V_1)
	\otimes GL(V_2)$ where the action of $(\rho_1\otimes\rho_2)(g)$
	on a basis element is
	\[(\rho_1\otimes\rho_2)(g) (e_i\otimes f_j) =
	(\rho_1(g)e_i)\otimes(\rho_2(g)f_j).\]
\end{definition}

\begin{prop}
	The character $\chi_1\otimes\chi_2$ of the representation
	$\rho_1\otimes\rho_2$ is just $\chi_1\cdot\chi_2$.
\end{prop}

\begin{proof}
	We need to calculate
	\[\chi_1\otimes\chi_2(g) =\sum_{i,j}
	\langle \rho_1(g)e_i\otimes\rho_2(g)f_j, e_i\otimes f_j\rangle,\]
	and the idea is that we can split up the summand as
	\[\langle \rho_1(g)e_i\otimes\rho_2(g)f_j, e_i\otimes f_j\rangle
	= \langle \rho_1(g)e_i, e_i\rangle\cdot\langle\rho_2(g)f_j,f_j\rangle.\]
	Apparently you prove this by looking at the matrix representations
	\todo{why?}
\end{proof}

\begin{definition}
	Given a unitary representation $\rho$ with matrix $R$,
	the \vocab{contragredient} representation $\rho^\ast$ has matrix
	$\ol R$ (complex conjugate all the entries).
\end{definition}

We can also take multiple tensor products; this gives the tensor power
representation.

\subsection{Tensor interlude: Alternating and symmetric tensors}
Given a vector space $V$, the alternating tensor is
\[\wedge^k V = \opname{span}\{e_{i_1}\wedge\cdots\wedge e_{i_k}\}_{i_1<\cdots<i_k},\]
where the wedge is defined as
\[x_1\wedge\cdots\wedge x_k = \sum_{\sigma\in \opname{Sym}(\{i_1,\dots,i_k\})}
\opname{sgn}(\sigma) x_{\sigma(1)}\otimes\cdots\otimes x_{\sigma(k)}.\]
The set $\opname{Sym}(S)$ is the set of all permutations of $S$.
The symmetric tensor $\vee^k V$ is the same thing but without
the $\opname{sgn}(\sigma)$ multiplier in the summand.

We have
\[V\tensor V = \wedge^2 V\oplus \vee^2 V,\]
noting that
\[e_i\tensor e_j = \half(e_i\tensor e_j + e_j\tensor e_i) +
\half(e_i\tensor e_j - e_j\tensor e_i)\]

In general the representations $\rho\tensor \rho$ splits
into invariant subspaces $\rho\wedge\rho\oplus\rho\vee\rho$,
and people study the symmetric and alternating power representations.

\subsection{Back to rep theory}
Recall that if $\rho$, $\rho'$ are representations of $G$ over spaces
$V$, $V'$ respectively, then given a $T\colon V\to V'$,
we can ``average'' out a homomorphism between the representations:
\[\widetilde T = \frac 1{|G|} \sum_{g\in G} (\rho'_g)\inv T\rho_g\]

\begin{exercise}
	What does this average look like if $V = V'$ and $\rho = \rho'$?
\end{exercise}

ok i died

------
\pagebreak
\setcounter{section}{5}
\section{Week 6 lmao}

\begin{theorem}
	Let $L/\QQ$ be a finite extension and let $B$ be
	the integral closure of $A$ in $L$.
	Then, in $B$, ideals have a unique factorization into prime ideals.
\end{theorem}

\begin{itemize}
	\ii Module -- ``vector space over a ring''
	\ii Noetherian module -- all submodules are of finite type,
	i.e. spanned by finitely many elements.
\end{itemize}

``Last time'' we proved that if $A$ is a ring and $E$ a module over $A$,
then $E$ is Noetherian iff $E'$ and $E/E'$ are Noetherian
for some submodule $E'\subseteq E$.
\begin{remark}
	A version of this theorem is in Artin for quadratic fields.
\end{remark}

\begin{corollary}
	Let $A$ be a ring and $E_1$, \dots, $E_n$ Noetherian $A$-modules.
	Then $E_1\times\cdots\times E_n$ is also Noetherian.
\end{corollary}

\begin{proof}
	For $n = 2$, we have $E_1\times E_2$ is Noetherian by taking
	$E' = E_1\times\{0\}\cong E_1$,
	then $E'\cong E_1$ and $E/E'\cong E_2$ are Noetherian.
	For bigger $n$ we induct.
\end{proof}

\begin{corollary}
	Let $A$ be a Noetherian ring and $E$ and $A$-module of finite type.
	Then $E$ is Noetherian.
\end{corollary}

\begin{proof}
	Let $y_1$, \dots, $y_n$ span $E$ so that
	$E = Ay_1 + \cdots + Ay_n$. There is a map $\phi\colon A^n\to E$ given by
	\[\phi(a_1, \dots, a_n) = a_1y_1 + \cdots + a_ny_n\]
	that is surjective. Let $\ker\phi = R$
	(called the \vocab{module of relations}); then $E\cong A^n/R$.
	$A$ is Noetherian, hence $A^n$ is Noetherian, hence
	$A^n/R\cong E$ is Noetherian, as desired.
\end{proof}

\begin{prop}
	Let $A$ be a Noetherian integrally closed ring and $K$
	its field of fractions. Let $L/K$ be a finite extension
	and let $B$ be the integral closure of $A$ in $L$.
	Assuming $K$ has characteristic zero, $B$ is an $A$-module of finite type
	and a Noetherian ring.
\end{prop}

\begin{proof}
	$B$ is a submodule of a free $A$-module.\todo{hm?}
	This implies that $B$ is finite type since $A$ is Noetherian.
	The ideals of $B$ are all $A$-submodules of $B$,
	so any family of ideals has a maximal element
	because $B$ is a Noetherian $A$-module, hence a Noetherian ring.
	This means the rings of integers of number fields are all Noetherian rings.
\end{proof}

Recall that an ideal $\mathfrak p$ of a ring $A$ is \vocab{prime} if
$xy\in\mathfrak p$ implies $x\in\mathfrak p$ or $y\in\mathfrak p$.
Inclusion is the ideal generalization of divisibility (cf. PIDs).
Recall also that $\mathfrak p$ is prime iff $A/\mathfrak p$ is integral.
Thus note that maximal ideals are obviously prime
(since quotienting by a maximal ideal yields a field).
One should think of the maximal ideals as analogous to irreducible elements.

\begin{lemma}
	Let $\mathfrak p$ be a prime ideal of a ring $A$ and let $A'\subseteq A$
	be a subring. Then $A'\cap\mathfrak p$ is a prime ideal of $A'$.
\end{lemma}
\begin{proof}
	Literally write down the definition:
	If $x,y\in A'$ and $xy\in A'\cap\mathfrak p$,
	then either $x\in A'\cap\mathfrak p$ or $y\in A'\cap \mathfrak p$.
\end{proof}

\begin{definition}
	If $\mathfrak a$, $\mathfrak b$ are ideals, then we define their product as
	\[\mathfrak a\cdot \mathfrak b :=\left\{\sum_{i=1}^k a_ib_i :
	a_i\in\mathfrak a,\;b_i\in\mathfrak b\right\}.\]
\end{definition}
\begin{exercise}
	Check that this is actually an ideal.
\end{exercise}

\begin{lemma}
	Let $\mathfrak p\subseteq A$ be an ideal and suppose that
	$\mathfrak p\supseteq \mathfrak a_1\mathfrak a_2\cdots\mathfrak a_n$
	where $\mathfrak a_1$, \dots, $\mathfrak a_n$ are ideals.
	Then $\mathfrak p\supseteq \mathfrak a_i$ for some $i$.
\end{lemma}
(Keep in mind the ``contains $\leftrightarrow$ divisibility'' conflation.)
\begin{proof}
	Suppose otherwise and let $a_i\in \mathfrak a_i\setdiff\mathfrak p$
	for each $i$. Then $a_1a_2\cdots a_n$ is in the product of the ideals,
	hence in $\mathfrak p$, by assumption, which is a contradiction.
\end{proof}

\begin{lemma}
	Let $A$ be a Noetherian integral domain.
	Every nonzero ideal of $A$ contains a product of prime ideals.
\end{lemma}
\begin{proof}
	Let $\Phi$ be the set of nonzero ideals not containing
	any product of prime ideals, and suppose for contradiction that
	this set is nonempty. Then the set has a maximal element $\mathfrak b$,
	which by the definition of $\Phi$ cannot be prime.
	Choose $x,y\in A\setdiff\mathfrak b$ so that $xy\in\mathfrak b$.
	Then $\mathfrak b + Ax$ and $\mathfrak b + Ay$ are ideals that are
	larger than $\mathfrak b$, so $\mathfrak b + Ax$ contains
	a product of prime ideals $\mathfrak p_1\cdots\mathfrak p_k$
	and $\mathfrak b + Ay$ contains a product of prime ideals
	$\mathfrak q_1\cdots\mathfrak q_\ell$. Then
	\[\mathfrak b\supseteq
	(\mathfrak b + Ax)(\mathfrak b + Ay)\supseteq
	\mathfrak p_1\cdots \mathfrak p_k\mathfrak q_1\cdots\mathfrak q_\ell,\]
	so $\mathfrak b$ contains a product of prime ideals.
\end{proof}

\begin{definition}
	Let $A$ be an integral domain and $K$ its field of fractions.
	An $A$-submodule $I$ of $K$ is called a \vocab{fractional ideal}
	if there exists $d\in A$ such that $dI\subseteq A$
	(so $d$ is a ``common denominator of $I$'').
\end{definition}

If $A$ is Noetherian, then all fractional ideals are finite type, since
$I\cong dI$. As a module, $dI\subseteq A$.

\begin{example}
	Some examples of fractional ideals:
	\begin{itemize}
		\ii Note that ideals of $A$ are fractional ideals with $d = 1$.
	
		\ii A submodule of $K$ of finite type is fractional, as,
		given generators $y_1$, \dots, $y_m$, one can clear their denominators.
	
		\ii If $I$, $I'$ are fractional ideals, then so are
		$I+I'$, $I\cap I'$, and $I\cdot I'$, where $I+I'$ and $I\cdot I'$
		are defined analogously to vanilla ideals.
	\end{itemize}
\end{example}

\begin{definition}
	Let $A$ be an integral domain. $A$ is a \vocab{Dedekind ring}
	if it is Noetherian, integrally closed, and all nonzero prime ideals
	are maximal. (For example, $\ZZ\subseteq\QQ$ is a Dedekind ring.)
\end{definition}

\begin{theorem}
	Let $A$ be a Dedekind ring, $K$ its field of fractions,
	$L$ a finite extension of $K$, $B$ the integral closure of $A$ in $L$.
	Suppose $K$ has characteristic zero. Then $B$ is a Dedekind ring
	and an $A$-module of finite type.
\end{theorem}

\begin{proof}
	We checked ``last class'' that the integral closure of $A$
	is integrally closed. We have already proven that $B$ is Noetherian
	and an $A$-module of finite type. To check the condition on prime ideals,
	let $\mathfrak p\ne 0$ be a prime ideal of $B$. Let $x$ be
	a nonzero element of $\mathfrak p$. Then $x$ is integral over $A$,
	so let $a_i\in A$ such that $x^n + \cdots + a_1x + a_0 = 0$
	is a polynomial of minimial degree of which $x$ is a root.
	Then $-x(x^{n-1} + a_{n-1}x^{n-2} + \cdots + a_1) = a_0\ne 0$
	and $a_0\in Bx\subseteq\mathfrak p$, hence $a_0\in\mathfrak p\cap A$,
	a nonzero prime ideal of $A$. This must be maximal because $A$ is Dedekind.
	
	Let $\phi\colon A\to B/\mathfrak p$ be the map $x\mapsto x+\mathfrak p$,
	whose kernel is $A\cap\mathfrak p$.
	Thus $A/A\cap\mathfrak p\hookleftarrow B/\mathfrak p$ is a field
	that is a subring of $B/\mathfrak p$. Since $B/\mathfrak p$ is integral
	over $A/A\cap\mathfrak p$, any $b\in B$ is a root
	of a monic polynomial with coefficients in $A$, which means that
	the projection $\ol b$ of $b$ to $B/\mathfrak p$ can also be written as
	a monic polynomial with coefficients in $A/A\cap\mathfrak p$.
	If we say that
	\[\ol {b^m} + \ol {a_{m-1}}{\ol b^{m-1}} + \cdots + \ol {a_0} = 0\]
	with $\ol{a_0}\ne\ol 0$, then we can subtract $\ol{a_0}$,
	factor out $\ol b$, and divide the other factor by $\ol{-a_0}$
	to procure a multiplicative inverse of $\ol b$, hence $B/\mathfrak p$
	is a field.
\end{proof}

\begin{theorem}
	Let $A$ be a Dedekind ring that is not a field.
	Then, every maximal ideal is invertible in fractional ideals;
	that is, if $\mathfrak m$ is maximal, there exists
	a fractional $\mathfrak m'$ such that $\mathfrak m\cdot\mf m' = A$.
\end{theorem}

\begin{proof}
	Let $\mf m\subset A$ be a nonzero maximal ideal
	(which exists because $A$ isn't a field).
	Let $\mf m'=\{x\in K : x\mf m\subseteq A\}$.
	We can check that $\mf m'$ is an $A$-submodule of $K$.
	Any $y\in\mf m$ serves as a denominator, so $\mf m'$ is a fractional ideal.
	
	We need to show that $\mf m\cdot\mf m' = A$. We know that
	$\mf m\cdot\mf m'$ is a subset of $A$ and, since $\mf m$ is an ideal,
	$\mf m = A\mf m\subseteq \mf m'\mf m$, so
	$\mf m\subseteq\mf m'\mf m\subseteq A$.
	Since $m$ is maximal, this implies that $\mf m'\mf m$ must be
	either $\mf m$ or $A$.
	
	If $\mf m\mf m' = \mf m$, then if $x\in\mf m'$, $x\mf m\subseteq \mf m$
	and in general $x^n\mf m\subseteq \mf m$ for all $n$.
	This means any $d\in M$ is a common denominator for the set of all $x^n$,
	so $A[x]$ is a fractional ideal. Since $A$ is Noetherian,
	$A[x]$ is an $A$-module of finite type, implying that
	$x$ is integral over $A$, implying $x\in A$.
	Thus $\mf m\mf m' = \mf m$ implies that $\mf m' = A$.
	
	We now need to rule of $\mf m' = A$.
	Let $a\in \mf m$. Then $A\mf a\subseteq\mf m$ contains a product of primes
	$\mf p_1\cdots\mf p_n$, so $\mf m$ contains one of the $\mf p_i$
	(WLOG $\mf p_1$). But $\mf p_1$ is prime hence maximal,
	so $\mf m = \mf p_1$. Let $\mf b = \mf p_2\cdots\mf p_n$.
	Then $Aa\supseteq\mf{mb}$ and $Aa$ doesn't contain $\mf b$.\todo{why}
	There exists $b\in\mathfrak b$ such that
	$b\notin Aa$, $\mf{m}b\subseteq Aa$, implying $ba\inv\in\mf m'$
	but $ba\inv\notin A$, so $\mf m'\ne A$.
\end{proof}

\begin{theorem}
	Let $A$ be a Dedekind domain.
	Every nonzero fractional ideal $\mf b$ has a unique expression
	\[\mf b = \prod_{\mf p\text{ maximal}}\mf p^{\nu_{\mf p}(\mf b)}.\]
	with only finitely many of the $\nu_{\mf p}(\mf b)$ being nonzero.
\end{theorem}

\begin{proof}
	For existence, pick $d\in A\setdiff\{0\}$.
	Then $d\cdot\mf b\subseteq A$. We have that $\mf b = (d\mf b)\cdot(dA)\inv$
	so it is sufficient to prove the theorem for ideals.
	Let $\Phi$ be the collection of nonzero ideals that are not
	products of prime ideals and suppose for contradiction that
	$\Phi$ is nonempty.
	Pick a maximal element of $\mf a\in\Phi$ (not maximal ideal).
	Then $\mf a$ is properly contained in some maximal ideal $\mf p$.
	Let $\mf p' = \mf p\inv$. We have
	\[\mf a\subseteq\mf p\implies \mf {ap}'\subseteq \mf{pp}' = A\]
	and also that $\mf p'$ properly contains $A$ and $\mf ap'\supseteq\mf a$.
	
	\begin{claim*}
		$\mf {ap}'\ne\mf a$.
	\end{claim*}
	\begin{proof}
		Otherwise, if $x\in\mf p'$, $x\mf a\subseteq \mf a$ implies that
		$x^n\mf a\subseteq\mf a$, which implies that $x$ is integral over $A$
		hence $x\in A$.
	\end{proof}
	Not all elements of $\mf p'$ are in $A$,
	thus $\mf {ap}'\notin\Phi$ by maximality.
	Thus $\mf{ap}'$ is a product of primes, and then we multiply by $\mf p$
	on the right to get $\mf a$ as a product of primes.
\end{proof}





















\end{document}